% Generated draft: requires mathematical and source review.
\documentclass[11pt]{article}
\usepackage[margin=25mm]{geometry}
\usepackage{fontspec}
\setmainfont{Latin Modern Roman}
\usepackage{amsmath,amssymb,mathtools,unicode-math}
\setmathfont{Latin Modern Math}
\usepackage{enumitem,needspace,xurl,hyperref}
\setlength{\parindent}{0pt}
\setlength{\parskip}{5pt}
\setlength{\emergencystretch}{3em}
\newcommand{\N}{\mathbb{N}}
\newcommand{\Z}{\mathbb{Z}}
\newcommand{\Q}{\mathbb{Q}}
\newcommand{\R}{\mathbb{R}}
\newcommand{\C}{\mathbb{C}}
\newcommand{\F}{\mathbb{F}}
\newcommand{\A}{\mathbb{A}}
\newcommand{\PP}{\mathbb{P}}
\newcommand{\E}{\mathbb{E}}
\newcommand{\eps}{\varepsilon}
\newcommand{\dd}{\,\mathrm{d}}
\begin{document}
\section*{UnsolvedMath ID: 30000168}
\textbf{Source problem code:} catalogue entry
\subsection{Definitions and mathematical statement}
Let \(K,N,M\), a map \(\tau:K\to N\), and a representation \(\sigma:K\to M\) be the objects intended in the source. The extension problem asks for conditions equivalent to
\[
\exists\rho:N\to M\text{ an admissible global representation}
\quad\text{such that}\quad
\rho(\tau(x))=\sigma(x)\quad\text{for every }x\in K.
\]
Thus ``\(\sigma\) is extensible to \(N\) via \(\tau\)'' means that this factorization exists. If the relevant representation classes admit composition, define
\[
\operatorname{res}_{\tau}:\operatorname{Rep}(N,M)\longrightarrow\operatorname{Rep}(K,M),
\qquad\rho\longmapsto\rho\circ\tau.
\]
The question is to characterize when \(\sigma\) lies in its image. This notation does not assume that representation means a group homomorphism into a matrix group; in a lattice workshop it could instead mean representation of forms or lattices by structure-preserving maps. The underlying category and the meaning of representation are not specified in the record.

The original passage additionally defines extension at a prime \(p\), referring to maps \(\rho^{(p)}:N_p\to M_p\) and a missing commutative diagram. If source-defined localizations and induced maps \(\tau_p,\sigma_p\) are available, the local condition would be \(\rho^{(p)}\circ\tau_p=\sigma_p\). The record does not define \(N_p,M_p,K_p\), identify the relevant primes, or specify whether completions or localizations are intended. It does not assert that local extensions at all primes imply a global extension; such a local-to-global principle would be an additional theorem to formulate and test, not a consequence of the definition.

The supplied literature triage checked 2026-08-21 says that the ambient objects, assumptions, and target theorem are too underspecified to identify a unique literature problem. Accordingly, the factorization question and its local terminology are preserved, without inventing hypotheses that make it a particular representation-extension conjecture.
\subsection{Short English statement}
Characterize when the given representation factors through \(\tau\) as \(\sigma=\rho\circ\tau\) for a global representation on \(N\). The structures being represented and the local-at-\(p\) conventions are absent, so no definite local-to-global theorem can yet be stated.
\subsection{Sources}
\begin{itemize}
\item Local UnsolvedMath registry, permanent ID 30000168, catalogue entry, \emph{Extension of Representations along \(\tau\)}.
\item \emph{Gitter und Anwendungen} (2005), Oberwolfach Reports, pp. 5--67, DOI \url{https://doi.org/10.4171/owr/2005/01}.
\item The supplied original passage includes the global and at-\(p\) extension definitions and refers to an absent Diagram (*). No separate authoritative theorem citation is supplied.
\item Supplied literature triage checked 2026-08-21 identifies the incomplete formulation; the corrected statement retains the explicit factorization question.
\end{itemize}
\end{document}
